\documentclass[11pt]{article}

%% Page geometry
\usepackage[a4paper, margin=1in]{geometry}

%% Core math + fonts
\usepackage{amsthm,amsmath,amsfonts,amssymb}
\usepackage{microtype}

%% Citations
\usepackage[authoryear, round]{natbib}

%% Hyperlinks
\usepackage[colorlinks, citecolor=blue, urlcolor=blue, linkcolor=blue]{hyperref}
\usepackage{xurl}

%% Graphics & figures
\usepackage{graphicx}
\graphicspath{{./images/}}
\usepackage{float}
\usepackage{subcaption}
\usepackage{booktabs}

%% Algorithms
\usepackage[linesnumbered, ruled, vlined]{algorithm2e}

%% TikZ for diagrams
\usepackage{tikz}
\usetikzlibrary{shapes,arrows,fit,calc,positioning}
\tikzset{box/.style={draw, diamond, thick, text centered, minimum height=0.5cm, minimum width=1cm}}
\tikzset{line/.style={draw, thick, -latex'}}

%% Misc
\usepackage{comment}
\usepackage{dsfont}

%% Theorem environments
\theoremstyle{plain}
\newtheorem{axiom}{Axiom}
\newtheorem{claim}[axiom]{Claim}
\newtheorem{theorem}{Theorem}[section]
\newtheorem{lemma}[theorem]{Lemma}

\theoremstyle{definition}
\newtheorem{definition}[theorem]{Definition}
\newtheorem*{example}{Example}
\newtheorem*{fact}{Fact}

\theoremstyle{remark}
\newtheorem{case}{Case}

%% Custom math commands
\newcommand{\E}{\mathbb{E}}
\newcommand{\1}{\mathds{1}}
\newcommand{\p}{\mathbb{P}}
\renewcommand{\H}{\mathcal{H}}
\newcommand{\T}{\mathcal{T}}
\newcommand{\ind}{\perp\!\!\!\perp}

%% File-specific additions
\usepackage{enumitem}

\newcommand{\X}{X}
\newcommand{\Y}{Y}
\newcommand{\A}{A}
\newcommand{\Sx}{S}
\newcommand{\U}{U}
\newcommand{\tauf}{\tau}
\newcommand{\cf}{c}
\newcommand{\gf}{g}
\newcommand{\muf}{\mu}
\newcommand{\mzero}{m_{0}}
\newcommand{\Spot}{S}
\newcommand{\RCT}{\mathrm{RCT}}
\newcommand{\RWD}{\mathrm{RWD}}
\newcommand{\RMST}{\mathrm{RMST}}

\title{Data fusion for heterogeneous survival effects:\\
goals, role of the HDP error, and positioning}
\author{Tijn Jacobs}
\date{\today}

\setlength{\parindent}{0pt}
\setlength{\parskip}{0.8em}


\begin{document}
\maketitle

\section*{Scope}

These notes accompany the proposed reframing of the project around two
explicit goals:
\begin{enumerate}[label=(G\arabic*),leftmargin=*]
  \item efficiency gain for heterogeneous treatment effect (HTE) estimation by
        combining a randomised trial (RCT, $\Sx=1$) with an observational
        study (RWD, $\Sx=0$);
  \item identifiable extrapolation of causal survival contrasts beyond the
        RCT follow-up horizon $t_{\RCT}$, using the longer follow-up of the RWD.
\end{enumerate}
Both goals are pursued under the AFT decomposition
\begin{equation}
  Y_i = \mzero(\X_i, \Sx_i) + \tauf(\X_i)\,\A_i
        + (1-\Sx_i)\,\A_i\,\cf(\X_i) + \sigma\varepsilon_i,
  \qquad Y_i = \log T_i,
  \label{eq:decomp}
\end{equation}
with $\mzero(\X,\Sx) = \muf(\X) + (1-\Sx)\gf(\X)$ in the four-forest version,
and $\varepsilon \sim P$ with $P$ assigned a hierarchical Dirichlet process
(HDP) prior across data sources.

\section{Pros of the reframing}

\paragraph{Pro 1: a clinically named problem, not just a methodological one.}
``Better HTE precision'' is a methodologist's pitch. ``Answering the long-term
survival question that the trial cannot answer alone'' is a clinical pitch
that names a concrete gap routinely encountered in health technology
assessment (HTA), oncology drug evaluation, and policy decision-making. The
gap is well-established: pivotal RCTs typically report 2--5 years of
follow-up, whereas registry data (NCDB, SEER, CPRD, claims databases) carry
10--20+ years \citep{Latimer2013, Jackson2023, Guyot2017}.

\paragraph{Pro 2: the AFT shift property yields a natural division of labour.}
On the log-time scale, $\tauf(\X)$ enters \eqref{eq:decomp} as a constant
covariate-specific shift. Once identified from short follow-up, the same
shift applies for all $t$. This makes the two information channels conceptually
orthogonal:
\begin{itemize}[leftmargin=*]
  \item the RCT identifies the \emph{shift} $\tauf(\X)$ (and protects
        $\tauf$ against confounding via the identification result;
        see Proposition~3 of the framework draft);
  \item the RWD identifies the \emph{shape} of the log-time distribution at
        long $t$ (through the error CDF $\Phi$) and the baseline $\mzero$.
\end{itemize}
The two sources therefore contribute complementary information, and not
merely additional sample size.

\paragraph{Pro 3: the HDP is no longer a robustness garnish; it carries
extrapolation.} Survival extrapolation is exquisitely sensitive to the
assumed error distribution. In a parametric AFT framework, lognormal,
Weibull, log-logistic, Gompertz, and generalised gamma families fit
short-term data equivalently well but yield wildly different tail behaviour,
and consequently wildly different RMST values evaluated past the trial
horizon. The NICE DSU TSD~14 workflow exists precisely because picking a tail
family by AIC on RCT data alone is unsafe \citep{Latimer2013, BellGorrod2019}.
The HDP prior on $\varepsilon$ \emph{is the mechanism} by which RWD follow-up
informs the tail: the shared atoms accumulate mass where both sources have
data; RWD-only mass shapes the tail; the hierarchical structure prevents
source-specific drift from contaminating the shared body of the error
distribution.

\paragraph{Pro 4: principled uncertainty quantification.} Long-term
extrapolations from parametric models often produce point estimates with
spuriously narrow credible intervals, because the family is treated as known.
A nonparametric error distribution propagates the genuine uncertainty about
tail shape into the posterior of $\Delta_{\RMST}(t^{*},\X)$ for
$t^{*} > t_{\RCT}$. As \citet{Jackson2023} put it: long-term estimates should
only be confident if there are strong long-term data. The HDP makes this
property automatic.

\paragraph{Pro 5: the two goals share one estimation machinery.} The same
posterior over $(\muf, \gf, \tauf, \cf, P, \sigma)$ yields both
the HTE estimate and the extrapolated survival contrast. There is no
two-stage handoff, no plug-in step where uncertainty is silently dropped,
and the confounding correction is in place throughout.

\section{Challenges that need to be addressed}

\paragraph{Challenge 1: constancy of the AFT shift over time.}
Under \eqref{eq:decomp}, $\tauf(\X)$ is a single number per covariate value;
extrapolation assumes the shift identified from $[0, t_{\RCT}]$ applies on
$(t_{\RCT}, \infty)$. This is a strong structural assumption. Treatments
whose effect wanes (or grows) over time violate it, and the literature has
explicit modelling devices for this case (waning treatment effects in
\citet{Jackson2023}; time-varying hazard ratios more generally). Two paths
forward:
\begin{itemize}[leftmargin=*]
  \item[(a)] retain constant $\tauf(\X)$ and present this as a transparent
        assumption with sensitivity analyses;
  \item[(b)] generalise to $\tauf(\X, t)$ via, e.g., a piecewise or
        spline-coefficient extension. This is feasible but breaks the
        clean ``shift'' interpretation.
\end{itemize}

\paragraph{Challenge 2: constancy of the confounding function over time.}
The decomposition \eqref{eq:decomp} treats $\cf(\X)$ analogously to
$\tauf(\X)$: a covariate-specific log-time shift that does not vary with $t$.
If selection bias has a different time signature than the treatment effect
itself --- e.g., frailty-driven early mortality differs across treatment
arms in the RWD but long-term survival does not --- a single $\cf(\X)$ is
mis-specified. Worth flagging, with sensitivity to time-varying confounding.

\paragraph{Challenge 3: period and calendar-time effects in the RWD.}
Long RWD follow-up necessarily mixes historical cohorts treated under earlier
standards of care. The patients contributing to the RWD tail are
\emph{older calendar} patients, not just \emph{older follow-up time}.
Without adjustment, this can violate Assumption~4 (transportability) for
reasons unrelated to unmeasured confounding. Three options:
condition on calendar year as a covariate; restrict the RWD window;
or model an explicit period effect alongside $\gf$.

\paragraph{Challenge 4: informative censoring in the RWD.}
Trial censoring is largely administrative and credibly non-informative;
RWD censoring includes loss-to-follow-up that may correlate with
prognosis. The HDP error does not fix this --- it assumes the residuals you
observe at long times are exchangeable within source. A sensitivity analysis
under a known censoring mechanism (or inverse-probability-of-censoring
weighting) seems necessary.

\paragraph{Challenge 5: who anchors the shared atoms?}
With a fully pooled DP, the RWD would dominate the error distribution
(because it is typically much larger and runs longer). With fully separate
DPs, the borrowing that gives extrapolation its power is lost. The
\emph{hierarchical} DP threads this needle, but the concentration parameters
governing how aggressively the source-specific distributions deviate from
the shared base must be either elicited or estimated. This is the
extrapolation analogue of the commensurate-prior shrinkage on $\gf$
controlled by $k_g$.

\paragraph{Challenge 6: identification of $\cf$ relies on the RCT.}
From Proposition~3(iii) of the framework draft, the RWD alone identifies
$\tauf + \cf$, not the two separately; the separation is anchored by the
RCT. With a small RCT and a large RWD, the prior on $\cf$ does substantial
work. Combined with extrapolation, this means: the RCT is doing
\emph{two} jobs simultaneously --- anchoring the shift and disambiguating
$\tauf$ from $\cf$. A small RCT may not be able to do both well. Worth
exploring via simulations that vary $n_{\RCT}/n_{\RWD}$.

\paragraph{Challenge 7: covariate overlap at long times.}
The HTE $\tauf(\X=x)$ is informative only at covariate values $x$ supported
by both sources. For extrapolation, an additional implicit overlap is
required: the RWD must contain survivors at long $t$ \emph{with covariate
profiles matching the RCT}. If the RWD long-followup is dominated by, say,
elderly patients who were not represented in the trial, the extrapolation
for the trial-representative population is effectively pure prior.

\paragraph{Challenge 8: connecting target-population definitions across goals.}
Under (G1), the natural target is the RWD population (representativeness
gain). Under (G2), the target is whichever population the HTA decision
concerns --- often the trial population, projected forwards in time. The
decomposition supports both, but the paper should be explicit about which
target population is used in each numerical result, since the survival
estimands in equations (23)--(26) of the framework draft contain $\cf$
only when the target is the RWD.

\section{Possible extensions}

\begin{enumerate}[label=(E\arabic*), leftmargin=*]
  \item \textbf{Time-varying treatment effect on the log-time scale.}
        Replace $\tauf(\X)$ by $\tauf(\X, t)$ via a spline or piecewise
        structure. This is the AFT-side analogue of waning treatment-effect
        models \citep{Jackson2023}.
  \item \textbf{Cure-fraction extension.}
        For oncology and immune-therapy settings, a mixture cure component
        is biologically motivated and is now standard in HTA extrapolation
        practice. Combining cure modelling with the four-forest decomposition
        is feasible and would broaden applicability.
  \item \textbf{Relative-survival / additive-hazards formulation.}
        Anchor the long-term hazard to known population mortality (life
        tables), so RWD-tail information competes against a hard external
        constraint rather than against the prior alone. Particularly natural
        for older populations with substantial competing mortality.
  \item \textbf{Sensitivity analysis on unmeasured confounding.}
        Allow $\cf(\X)$ to be informed by a prior derived from negative
        controls, or perform a tipping-point analysis over the magnitude of
        $\cf$ consistent with observed RCT-RWD discrepancy.
  \item \textbf{Multiple data sources.}
        Generalise from $\Sx \in \{0,1\}$ to a categorical source indicator,
        with the HDP base distribution shared across all sources and
        source-specific deviations regulated hierarchically. Several
        registries plus a trial; or a trial, an EHR cohort, and a claims
        database.
  \item \textbf{Aggregate external evidence.}
        Allow inclusion of summary survival data (e.g., from published KM
        curves reconstructed via \citet{Guyot2012}) alongside individual-level
        data, in the spirit of survextrap. Useful when long-term IPD are
        unavailable but published curves exist.
  \item \textbf{Posterior projection summaries for the HTE.}
        Map the high-dimensional posterior over $\tauf(\X)$ into an
        interpretable summary (a CART tree, a small set of subgroup rules)
        via the projection approach of \citet{Woody2020}, and report
        how those subgroup summaries shift when the RWD is included.
        This directly answers ``how did borrowing change the
        heterogeneity story?''
  \item \textbf{Calendar-time covariate.}
        Adjust for cohort year inside the BART forests, mitigating
        Challenge~3.
  \item \textbf{Decision-analytic outputs.}
        Pipe the posterior of $\Delta_{\RMST}(t^*, \X)$ into life-years
        gained and (with utility weights) QALY contrasts. This is what HTA
        agencies actually want; the method already produces it.
\end{enumerate}

\section{Related work}

The reframing places the project at the intersection of three literatures
that have historically not communicated very much: causal data fusion,
Bayesian survival extrapolation for HTA, and flexible Bayesian survival
modelling. We sketch each below and indicate where the proposed method sits.

\subsection{Survival extrapolation for HTA}

The foundational reference is NICE DSU Technical Support Document~14
\citep{Latimer2013}, which codified the practice of fitting standard
parametric AFT/PH families (exponential, Weibull, log-logistic, lognormal,
Gompertz, generalised gamma) to trial data and selecting among them by
fit and clinical plausibility. The TSD~14 process is widely followed but
also widely criticised: fits in the observed window are nearly
indistinguishable, while extrapolations diverge dramatically
\citep{BellGorrod2019, Bagust2014}.

External data have been used in HTA extrapolation in several distinct ways:
\begin{itemize}[leftmargin=*]
  \item \textbf{General-population mortality} as a lower bound on the
        hazard via additive-hazards (relative survival) models;
  \item \textbf{Registry conditional survival} from cancer registry
        databases used as external survival information for cancer
        \citep{Guyot2017};
  \item \textbf{Expert elicitation} treated as prior data on long-term hazard
        \citep{Guyot2017};
  \item \textbf{Longer-follow-up surrogate trials} such as the
        cilta-cel/LEGEND-2 example \citep{Vyas2023}, in which a phase~I
        study with longer follow-up informs extrapolation of a pivotal
        phase~II trial through informative priors on shape parameters.
\end{itemize}

The most directly comparable methodological piece is \citet{Jackson2023}'s
\texttt{survextrap}, which fits an M-spline hazard with a hierarchical prior
on the coefficients, jointly to individual-level RCT data and external
(possibly aggregate) data using Bayesian multi-parameter evidence synthesis.
\texttt{survextrap} supports proportional and flexible non-proportional
hazards, waning treatment effects, cure models, and additive-hazards
formulations \citep{Jackson2023}. Empirical evaluations
\citep{Timmins2025} report that the model gives accurate long-term
extrapolations when long-term external data on the patient population of
interest are included, and quantifies structural uncertainty appropriately
when they are not.

\paragraph{Where the proposed method differs from \texttt{survextrap}.}
\texttt{survextrap} does not address unmeasured confounding in the external
data; the external sources it accommodates are general-population mortality,
registry conditional survival, and elicited beliefs --- not a confounded
real-world data of the \emph{same} treatment. The proposed framework
adds an explicit confounding-function correction $\cf(\X)$, allowing the
external data to be a contemporaneous RWD subject to non-randomised
treatment assignment, not merely a background hazard reference. The
modelling style also differs: M-spline hazards versus BART forests on
AFT components, and a flexible parametric error versus an HDP nonparametric
error.

\subsection{Causal data fusion under unmeasured confounding}

The methodological core --- a confounding function $\cf(\X)$ correcting RWD
contrasts that do not hold unconfoundedness --- is established in
\citet{YangLiuZengWang2025}, who derive semiparametric efficiency bounds for
joint estimation of $\tauf$ and $\cf$. \citet{Kallus2018} propose the
``experimental grounding'' framework, in which an RCT supplies a bias
function correcting an observational HTE estimator. \citet{Yao2024} extend
data fusion to partial RWD treatment information. \citet{vanderLaan2024}
develop adaptive TMLE for RCT--RWD fusion that adapts borrowing to detected
bias.

The closest match to the proposed work is \citet{Ye2025}, who develop an
integrative AFT framework for high-dimensional RCT and RWD with censoring
and hidden confounding, using Stute-weighted penalised regression with
separate regularisation for $\tauf$ and the confounding component. The
proposed method is essentially a Bayesian-nonparametric counterpart to
\citet{Ye2025}: BART ensembles in place of penalised parametric models,
nonparametric error in place of a Gaussian or assumed AFT family, and
commensurate-prior structure on the baseline split between sources.

A practical caveat from \citet{Lin2024}: bias-correction methods for
marginal ATE estimation in this literature have not delivered
RMSE improvements over RCT-only analysis in their simulations, and the
methods' theoretical advantage is for the \emph{heterogeneous} effect
function rather than its average. The proposed framing targets HTE
plus long-term extrapolation, both of which are tasks for which RCT-only
inference is genuinely deficient --- the latter, simply because it cannot
be done from the trial alone.

\subsection{Bayesian flexible survival modelling}

The BART-based modelling of $\muf$, $\gf$, $\tauf$, $\cf$ rests on the
Bayesian causal forest (BCF) line of work \citep{hahn2020bayesian}, which introduced
the prognostic/treatment-effect decomposition with separate prior scales
that the four-forest version generalises. The HDP error builds on the
centred stick-breaking mixture of \citet{yang2010semiparametric} and the wider
Bayesian-nonparametric survival literature; the hierarchical extension
specifically lets the source label index the group hierarchy, with shared
atoms across sources.

\subsection{Positioning summary}

\begin{table}[tb]
\centering\small
\begin{tabular}{lccc}
\toprule
& Confounding & Long-term      & Nonparametric \\
& correction  & extrapolation  & error         \\
\midrule
\citet{Latimer2013} (TSD 14 workflow)        & --- & yes & ---  \\
\citet{Guyot2017}                            & --- & yes & ---  \\
\citet{Jackson2023} (\texttt{survextrap})    & --- & yes & flexible \\
\citet{Ye2025}                               & yes & limited$^\dagger$ & --- \\
\citet{YangLiuZengWang2025}                  & yes & ---       & --- \\
\citet{Kallus2018}, \citet{Yao2024}          & yes & ---       & --- \\
\textbf{This work}                           & \textbf{yes} & \textbf{yes} & \textbf{yes (HDP)} \\
\bottomrule
\end{tabular}
\caption{Schematic positioning. $\dagger$: \citet{Ye2025} targets HTE
estimation in the joint AFT framework; long-term extrapolation is not its
stated goal but is in principle compatible with the AFT structure.}
\end{table}

The niche is clear: the HTA-pragmatic literature handles extrapolation
flexibly but does not correct for confounding in the external source; the
causal-fusion literature corrects for confounding but does not target
extrapolation; and few methods combine the two with a fully nonparametric
error distribution. The proposed method sits deliberately in that gap.

\bibliographystyle{apalike}
\bibliography{../general/references}

\end{document}
